% ============================================================
% 1.4. Постановка задачи
% ============================================================
\subsection{Постановка задачи}

\subsubsection{Математическая модель плоского изгиба балки}

Рассматривается задача плоского поперечного изгиба прямолинейной призматической балки длиной~$L$ с постоянной изгибной жёсткостью $EI$~\cite{feodosev, timoshenko}. Балка нагружена совокупностью внешних силовых факторов: сосредоточенных сил $F_i$ в точках $a_i$, распределённых нагрузок $q_j(x)$ на участках $[x_{j1}, x_{j2}]$ и сосредоточенных моментов $M_{0k}$ в точках $b_k$.

\textbf{Вывод основного дифференциального уравнения изогнутой оси балки.} Получим уравнение, связывающее прогиб $v(x)$ с изгибающим моментом $M(x)$, из геометрических и физических предпосылок~\cite{feodosev, timoshenko, gere2018}.

\textit{Геометрическое соотношение.} Кривизна изогнутой оси балки в точке с координатой $x$ выражается через прогиб $v(x)$ по формуле:
\begin{equation}
  \kappa(x) = \frac{v''(x)}{\left[1 + (v'(x))^2\right]^{3/2}}.
  \label{eq:kappa_full}
\end{equation}
В предположении малых прогибов ($(v')^2 \ll 1$), соответствующем допущению о~малых перемещениях, знаменатель можно заменить единицей, и кривизна упрощается:
\begin{equation}
  \kappa(x) \approx v''(x) = \frac{d^2 v}{dx^2}.
  \label{eq:kappa_approx}
\end{equation}

\textit{Физико-геометрическое соотношение.} Ранее было получено соотношение~(\ref{eq:kappa_M}), связывающее кривизну с изгибающим моментом через жёсткость~$EI$:
\[
  \kappa(x) = \frac{M(x)}{E I}.
\]

\textit{Приравнивание и итоговое уравнение.} Приравнивая~(\ref{eq:kappa_approx}) и~(\ref{eq:kappa_M}), получаем основное дифференциальное уравнение изогнутой оси балки (уравнение Эйлера--Бернулли)~\cite{timoshenko1983, gere2018}:
\begin{equation}
  E I \, \frac{d^2 v}{dx^2} = M(x),
  \label{eq:euler_bernoulli}
\end{equation}
\begin{conditions}
  E & модуль упругости материала, \textit{Па}; \\
  I & осевой момент инерции поперечного сечения, \textit{м}$^4$; \\
  v(x) & прогиб (вертикальное перемещение точек оси), \textit{м}; \\
  M(x) & изгибающий момент в сечении $x$, \textit{Н$\cdot$м}.
\end{conditions}

Полученное уравнение имеет ясный физический смысл: чем больше изгибающий момент $M(x)$, тем сильнее искривлена ось балки; чем больше жёсткость $EI$, тем меньше кривизна при том же моменте. Уравнение (\ref{eq:euler_bernoulli}) является основой дальнейшего численного интегрирования для определения прогибов и углов поворота.

\textbf{Альтернативные формы уравнения изогнутой оси.} Используя дифференциальные соотношения равновесия, уравнение (\ref{eq:euler_bernoulli}) можно записать также через поперечную силу и через распределённую нагрузку:
\begin{equation}
  E I \, \frac{d^3 v}{dx^3} = Q(x),
  \label{eq:euler_Q}
\end{equation}
\begin{equation}
  E I \, \frac{d^4 v}{dx^4} = -q(x).
  \label{eq:euler_q}
\end{equation}
Уравнение четвёртого порядка~(\ref{eq:euler_q}) удобно при известной функции распределённой нагрузки и сводит расчёт балки к последовательному интегрированию дифференциального уравнения с четырьмя постоянными интегрирования, определяемыми из граничных условий.

\textbf{Вывод дифференциальных соотношений равновесия.} Получим зависимости между интенсивностью распределённой нагрузки $q(x)$, поперечной силой $Q(x)$ и изгибающим моментом $M(x)$ из условия равновесия бесконечно малого элемента балки~\cite{feodosev, belyaev, darkov}.

Выделим из балки элемент длиной $dx$, расположенный между сечениями $x$ и $x + dx$. На левом сечении действуют поперечная сила $Q(x)$ и изгибающий момент $M(x)$, на правом сечении --- $Q(x + dx) = Q + dQ$ и $M(x + dx) = M + dM$. На элементе распределена нагрузка интенсивности $q(x)$, результирующая сила которой на участке $dx$ равна $q \cdot dx$.

\textbf{Уравнение равновесия сил} (сумма проекций на ось $y$):
\begin{equation}
  Q - (Q + dQ) - q \cdot dx = 0.
  \label{eq:eq_force_deriv}
\end{equation}

Раскрывая скобки и деля на $dx$, получаем связь между распределённой нагрузкой и поперечной силой:
\begin{equation}
  \frac{dQ}{dx} = -q(x).
  \label{eq:shear_load_relation}
\end{equation}

Физический смысл: производная поперечной силы по координате равна интенсивности распределённой нагрузки, взятой с обратным знаком. Отсюда следствия: на участке без распределённой нагрузки ($q = 0$) поперечная сила постоянна, эпюра $Q(x)$ --- горизонтальная линия; на участке с равномерной нагрузкой ($q = \text{const}$) эпюра $Q(x)$ --- прямая линия с наклоном $-q$; на участке с линейной нагрузкой эпюра $Q(x)$ --- парабола второго порядка~\cite{feodosev}.

\textbf{Уравнение равновесия моментов} (сумма моментов относительно правого сечения $x + dx$):
\begin{equation}
  M + Q \cdot dx - (M + dM) - q \cdot dx \cdot \frac{dx}{2} = 0.
  \label{eq:eq_moment_deriv}
\end{equation}

Пренебрегая слагаемым второго порядка малости $q \cdot (dx)^2 / 2$ и деля на $dx$, получаем связь между изгибающим моментом и поперечной силой:
\begin{equation}
  \frac{dM}{dx} = Q(x).
  \label{eq:moment_shear_relation}
\end{equation}

Физический смысл: производная изгибающего момента по координате равна поперечной силе. Отсюда вытекает фундаментальное следствие: в точках, где поперечная сила обращается в нуль ($Q(x_0) = 0$), производная $M'(x_0) = 0$, что соответствует экстремуму изгибающего момента.

\textbf{Вторая производная изгибающего момента.} Дифференцируя~(\ref{eq:moment_shear_relation}) и подставляя~(\ref{eq:shear_load_relation}), получаем:
\begin{equation}
  \frac{d^2 M}{dx^2} = \frac{dQ}{dx} = -q(x).
  \label{eq:moment_load_relation}
\end{equation}

Это соотношение связывает вторую производную изгибающего момента непосредственно с интенсивностью нагрузки и позволяет судить о типе эпюры $M(x)$ по характеру распределения нагрузки.
\begin{conditions}
  Q(x) & поперечная сила в сечении $x$, \textit{Н}; \\
  q(x) & интенсивность распределённой нагрузки, \textit{Н/м}.
\end{conditions}

Из соотношения~(\ref{eq:moment_shear_relation}) следует фундаментальное правило: \textit{экстремум изгибающего момента достигается в сечении, где поперечная сила обращается в нуль} ($Q(x_0) = 0 \Rightarrow M'(x_0) = 0$)~\cite{feodosev}.

\subsubsection{Уравнения равновесия}

Для статически определимой системы неизвестные опорные реакции определяются из уравнений статики~\cite{aleksandrov, belyaev}. В плоской задаче имеется три уравнения равновесия:
\begin{equation}
  \sum F_x = 0, \qquad \sum F_y = 0, \qquad \sum M_A = 0.
  \label{eq:equilibrium}
\end{equation}

Число неизвестных реакций для каждого типа балки:
\begin{itemize}
  \item \textbf{Шарнирно-опёртая балка:} $R_A$ и $R_B$ (2~неизвестных, при отсутствии горизонтальных нагрузок достаточно 2~уравнений $\sum F_y = 0$ и $\sum M_A = 0$).
  \item \textbf{Консольная балка:} $R_{\text{зад}}$ и $M_{\text{зад}}$ (2~неизвестных).
  \item \textbf{Балка с вылетом:} $R_A$ и $R_B$ (2~неизвестных, аналогично шарнирно-опёртой).
\end{itemize}

Во всех случаях система статически определима: число уравнений равновесия не менее числа неизвестных реакций~\cite{darkov}.

\subsubsection{Метод сечений для построения эпюр}

Построение эпюр внутренних силовых факторов осуществляется методом сечений~\cite{feodosev, pisarenko}. Метод основан на мысленном рассечении балки плоскостью, перпендикулярной оси, с последующим рассмотрением равновесия отделённой части (метод ROSE --- Resect, Omit, Substitute, Equilibrate)~\cite{hibbeler2016}.

Для произвольного сечения $x$ суммируются вклады всех нагрузок и реакций, расположенных \textit{левее} точки сечения (метод левой части):

\textbf{От сосредоточенной силы $F_i$ в точке $a_i$ (при $x \geq a_i$):}
\begin{equation}
  Q(x) \mathrel{-}= F_i, \qquad M(x) \mathrel{-}= F_i \cdot (x - a_i).
  \label{eq:section_point}
\end{equation}

\textbf{От равномерной нагрузки $q$ на участке $[x_1, x_2]$ (при $x > x_1$):}
\begin{equation}
  x_{\text{эфф}} = \min(x, x_2), \quad l_{\text{эфф}} = x_{\text{эфф}} - x_1,
  \label{eq:section_uniform_1}
\end{equation}
\begin{equation}
  \Delta F = q \cdot l_{\text{эфф}}, \quad x_c = x_1 + \frac{l_{\text{эфф}}}{2},
  \label{eq:section_uniform_2}
\end{equation}
\begin{equation}
  Q(x) \mathrel{-}= \Delta F, \qquad M(x) \mathrel{-}= \Delta F \cdot (x - x_c).
  \label{eq:section_uniform_3}
\end{equation}

\textbf{От треугольной нагрузки $q_1 \to q_2$ на участке $[x_1, x_2]$ (при $x > x_1$):}
\begin{equation}
  t = \frac{l_{\text{эфф}}}{l}, \quad q_x = q_1 + (q_2 - q_1) \cdot t,
  \label{eq:section_tri_1}
\end{equation}
\begin{equation}
  \Delta F = \frac{q_1 + q_x}{2} \cdot l_{\text{эфф}}, \quad x_c = x_1 + l_{\text{эфф}} \cdot \frac{q_1 + 2 q_x}{3(q_1 + q_x)}.
  \label{eq:section_tri_2}
\end{equation}

\textbf{От сосредоточенного момента $M_0$ в точке $b$ (при $x \geq b$):}
\begin{equation}
  M(x) \mathrel{-}= M_0.
  \label{eq:section_moment}
\end{equation}

\subsubsection{Определение деформаций}

Прогиб $v(x)$ и угол поворота $\theta(x)$ определяются двойным интегрированием уравнения~(\ref{eq:euler_bernoulli})~\cite{feodosev, popov1998}:
\begin{equation}
  \theta(x) = \frac{1}{EI} \int_0^x M(t) \, dt + C_1,
  \label{eq:theta}
\end{equation}
\begin{equation}
  v(x) = \int_0^x \theta(t) \, dt + C_2.
  \label{eq:deflection}
\end{equation}

Постоянные интегрирования $C_1$ и $C_2$ определяются из граничных условий~\cite{aleksandrov}:

\textbf{Консольная балка:} из условий $v(0) = 0$ и $\theta(0) = 0$:
\begin{equation}
  C_1 = 0, \qquad C_2 = 0.
  \label{eq:const_cant}
\end{equation}

\textbf{Шарнирно-опёртая балка} и \textbf{балка с вылетом:} из условий $v(0) = 0$ и $v(L) = 0$:
\begin{equation}
  C_2 = 0, \qquad C_1 = -\frac{v_{\text{raw}}(L)}{L},
  \label{eq:const_ss}
\end{equation}
где $v_{\text{raw}}(L)$ --- значение прогиба в точке $x = L$, вычисленное при $C_1 = 0$.

\subsubsection{Формулировка задачи в терминах вычислительной математики}

Задача формулируется как последовательность вычислительных этапов~\cite{bahvalov, samarsky}:

\begin{enumerate}
  \item \textbf{Дискретизация.} Ось балки разбивается на $N = 500$ равномерных точек:
  \begin{equation}
    x_k = k \cdot \frac{L_{\text{полн}}}{N - 1}, \quad k = 0, 1, \ldots, N-1,
    \label{eq:discretization}
  \end{equation}
  где $L_{\text{полн}} = L + c$ для балки с вылетом ($c$ --- длина вылета), $L_{\text{полн}} = L$ в остальных случаях.

  \item \textbf{Вычисление эпюр.} Для каждой точки $x_k$ рассчитываются значения $Q(x_k)$ и $M(x_k)$ методом суперпозиции (сложение вкладов всех нагрузок и реакций).

  \item \textbf{Численное интегрирование.} Угол поворота $\theta(x)$ и прогиб $v(x)$ вычисляются методом трапеций:
  \begin{equation}
    I_{k} = I_{k-1} + \frac{f(x_{k-1}) + f(x_k)}{2} \cdot (x_k - x_{k-1}).
    \label{eq:trapezoid}
  \end{equation}

\end{enumerate}

Точность аппроксимации определяется шагом дискретизации $h = L_{\text{полн}} / (N-1)$. При $N = 500$ и $L = 10$~м шаг составляет $h \approx 0{,}02$~м, что обеспечивает погрешность метода трапеций порядка $O(h^2)$~\cite{bahvalov}.

\subsubsection{Обоснование выбранной математической модели}

Сформулированная математическая модель основывается на следующих сведённых воедино предпосылках:

\textbf{1. Механические предпосылки.} Балка рассматривается как одномерный континуум (прямолинейный стержень) в условиях плоского изгиба. Распределение напряжений по сечению подчиняется формуле~(\ref{eq:sigma_M_y}) и линейно зависит от координаты $y$ относительно нейтральной оси. Сдвигами и инерционными эффектами пренебрегают.

\textbf{2. Геометрические предпосылки.} Деформации считаются малыми: $\varepsilon \ll 1$, $v \ll L$, $(v')^2 \ll 1$. Поперечные сечения остаются плоскими и перпендикулярными деформированной оси (гипотеза Бернулли). Это позволяет применять линеаризованное уравнение изогнутой оси~(\ref{eq:euler_bernoulli}).

\textbf{3. Физические предпосылки.} Материал считается однородным, изотропным и линейно-упругим, подчиняющимся закону Гука $\sigma = E \varepsilon$. Пластические деформации, ползучесть, термические напряжения не учитываются.

\textbf{4. Статическая определимость.} Рассматриваются только системы, в которых число неизвестных реакций равно числу уравнений равновесия. Это позволяет определять реакции из условий статики без привлечения уравнений совместности деформаций и не требует решения больших систем линейных уравнений.

\textbf{5. Возможность применения принципа суперпозиции.} Благодаря линейности дифференциального уравнения~(\ref{eq:euler_bernoulli}) и закона Гука, к задаче применим принцип суперпозиции. Это позволяет разложить сложную задачу на несколько простых: вычислить эпюры от каждой нагрузки отдельно и сложить результаты.

Указанные предпосылки определяют класс решаемых задач и одновременно --- границы применимости разрабатываемой платформы, изложенные в подразделе~1.2.

\subsubsection{Классификация задач расчёта балок}

В зависимости от того, какие величины известны и какие требуется определить, задачи расчёта балок делятся на несколько классов~\cite{feodosev, aleksandrov}:

\textbf{Прямая задача расчёта.} Известны: тип балки, геометрия и нагрузки. Требуется определить: реакции опор, эпюры $Q(x)$, $M(x)$, $N(x)$, прогибы $v(x)$ и углы поворота $\theta(x)$. Эта задача является основной для курса сопротивления материалов и реализована в платформе SopromatLab.

\textbf{Задача проверки прочности.} Дополнительно к прямой задаче задаются допускаемые напряжения $[\sigma]$. Требуется проверить, выполняется ли условие прочности~(\ref{eq:strength_condition}). Такая задача относится к дальнейшему развитию расчётного модуля.

\textbf{Задача подбора сечения.} Задаются максимальные нагрузки и допускаемые напряжения. Требуется подобрать момент сопротивления сечения $W$ и, соответственно, тип и размер профиля. Данная задача не входит в реализованный функционал платформы.

\textbf{Задача определения допускаемой нагрузки.} Задаются геометрия и допускаемые напряжения. Требуется определить максимальную допустимую величину нагрузки, при которой $\sigma_{\max} \leq [\sigma]$. Данная задача сводится к прямой задаче с вариативной величиной нагрузки и может быть рассмотрена как направление дальнейшего развития.

Все перечисленные классы задач сводятся к многократному выполнению прямой задачи расчёта, что соответствует интерактивному режиму работы платформы: студент может менять параметры модели и наблюдать мгновенное обновление результатов.
